% general-case-algorithm-proof.tex — \input'd from proofs.tex
% Conventions: thesis/CLAUDE.md
%
% Sources: Jörn's dictation + HK2017 paper. Notation per correspondence.tex.
%
% Structure: Proof of main theorem, supporting lemmas, algorithm correctness

% Jörn: text approved (692c271) — section intro
We prove Theorem~\ref{thm:main}: the algorithm in
Section~\ref{subsec:general-case-algorithm} correctly computes
$c_{\mathrm{EHZ}}(K)$.

% Jörn: text approved (692c271) — from \begin{theorem*} to \end{theorem*}
\begin{theorem*}[Theorem~\ref{thm:main}, restated]
Let $K \subset \mathbb{R}^4$ be a polytope
with $F$ facets, outward unit normals $n_1, \ldots, n_F \in S^3$,
and heights $h_1, \ldots, h_F > 0$.
Then
\begin{equation}\tag{\ref{eq:ehz-main}}
  c_{\mathrm{EHZ}}(K) = \frac{1}{2}
  \left[
    \max_{\substack{
      \sigma \in S_F \\
      \beta \in M(K)
    }}
    \sum_{1 \leq j < i \leq F}
      \beta_{\sigma(i)} \beta_{\sigma(j)} \,
      \omega_0(n_{\sigma(j)}, n_{\sigma(i)})
  \right]^{-1},
\end{equation}
where
\begin{equation}\tag{\ref{eq:constraint-set}}
  M(K) = \left\{
    (\beta_i)_{i=1}^{F} :
    \beta_i \geq 0, \quad
    \sum_{i=1}^{F} \beta_i h_i = 1, \quad
    \sum_{i=1}^{F} \beta_i n_i = 0
  \right\}.
\end{equation}
\end{theorem*}

\begin{proof}[Proof of Theorem~\ref{thm:main}]
% Jörn: text approved (b7a8bc0) — from \begin{proof} to \end{proof}
% QC: combines thm:simple-minimizer, lem:shoelace, and a variable
%   substitution. Each step cites specific results and states which
%   properties are used. The critical calculation is beta^T H beta = 1/T.
%   Verified against HK2017 proof of Theorem 1.1 (lines 455--478):
%   same substitution beta_{sigma(i)} = T_i / h_{sigma(i)}, same
%   identity, same conclusion.

\emph{Step~1: Reduce to simple orbits.}
By Theorem~\ref{thm:simple-minimizer},
$c_{\mathrm{EHZ}}(K) = \min\{ A(\gamma) :
\gamma \text{ is a simple Reeb orbit on } \partial K \}$.
A simple orbit is described by a tuple
$(S, \sigma, \tau, b)$ where $S \subseteq \{1, \ldots, F\}$ is
nonempty, $\sigma$ is a cyclic ordering of~$S$,
$\tau_1, \ldots, \tau_{|S|} > 0$ are positive durations, and
$b = \gamma(0) \in \partial K$
(Remark~\ref{rem:simple-orbit-data}).
By Lemma~\ref{lem:shoelace}, the action depends only on
$(\sigma, \tau)$, not on~$b$: closure requires
$\sum_{k=1}^{|S|} \tau_k R_{\sigma(k)} = 0$, and
$A(\gamma) = T = \sum_k \tau_k$
(Remark~\ref{rem:contact-normalization}).

\emph{Step~2: Apply the shoelace formula.}
By Definition~\ref{def:action} and $J_0^\top = -J_0$
(Remark~\ref{rem:j0-properties}),
$2A(\gamma) = \int_0^T \langle -J_0 \dot{\gamma},\,
\gamma \rangle\, dt$.
By Lemma~\ref{lem:shoelace}:
\begin{equation}\label{eq:shoelace-orbit}
  2T = \sum_{1 \leq j < i \leq |S|}
    \tau_j\, \tau_i\,
    \omega_0(R_{\sigma(j)},\, R_{\sigma(i)}).
\end{equation}

\emph{Step~3: Expand Reeb vectors.}
Since $R_i = \frac{2}{h_i} J_0 n_i$
(Definition~\ref{def:reeb-vectors}) and
$\omega_0(J_0 u, J_0 v) = \omega_0(u, v)$
(Remark~\ref{rem:j0-properties}):
\[
  \omega_0(R_i, R_j)
  = \frac{4}{h_i h_j}\, \omega_0(J_0 n_i,\, J_0 n_j)
  = \frac{4}{h_i h_j}\, \omega_0(n_i, n_j).
\]
Substituting into~\eqref{eq:shoelace-orbit}:
\begin{equation}\label{eq:shoelace-normals}
  2T = 4 \sum_{1 \leq j < i \leq |S|}
    \frac{\tau_j\, \tau_i}
         {h_{\sigma(j)}\, h_{\sigma(i)}}\,
    \omega_0(n_{\sigma(j)},\, n_{\sigma(i)}).
\end{equation}

\emph{Step~4: Variable substitution $\tau \to \beta$.}
Define
\begin{equation}\label{eq:beta-def}
  \beta_{\sigma(k)} \coloneqq
    \frac{\tau_k}{T\, h_{\sigma(k)}},
  \qquad k = 1, \ldots, |S|,
\end{equation}
and $\beta_i = 0$ for $i \notin S$.
We verify $\beta \in M(K)$
(the constraint set~\eqref{eq:constraint-set}):
\begin{enumerate}
\item \emph{Non-negativity:}
  $\beta_i \geq 0$ for all~$i$, with
  $\beta_{\sigma(k)} > 0$ for $k \in \{1,\ldots,|S|\}$.
\item \emph{Height normalization:}
  $\displaystyle\sum_{i=1}^F \beta_i h_i
  = \sum_{k=1}^{|S|}
    \frac{\tau_k}{T h_{\sigma(k)}} \cdot h_{\sigma(k)}
  = \frac{1}{T} \sum_k \tau_k = 1$.
\item \emph{Closure:}
  The orbit closure
  $\sum_k \tau_k R_{\sigma(k)} = 0$ gives
  $\sum_k \frac{2\tau_k}{h_{\sigma(k)}} J_0 n_{\sigma(k)} = 0$.
  Since $J_0$ is invertible:
  $\sum_k \frac{\tau_k}{h_{\sigma(k)}} n_{\sigma(k)} = 0$.
  Dividing by~$T$:
  $\sum_{i=1}^F \beta_i n_i = 0$.
\end{enumerate}

\emph{Step~5: Derive the action formula.}
Substitute $\tau_k = T h_{\sigma(k)} \beta_{\sigma(k)}$
into~\eqref{eq:shoelace-normals}:
\begin{align*}
  2T &= 4 \sum_{j < i}
    \frac{(T h_{\sigma(j)} \beta_{\sigma(j)})\,
          (T h_{\sigma(i)} \beta_{\sigma(i)})}
         {h_{\sigma(j)}\, h_{\sigma(i)}}\,
    \omega_0(n_{\sigma(j)},\, n_{\sigma(i)}) \\
  &= 4T^2 \sum_{1 \leq j < i \leq |S|}
    \beta_{\sigma(i)}\, \beta_{\sigma(j)}\,
    \omega_0(n_{\sigma(j)},\, n_{\sigma(i)}).
\end{align*}
The sum on the right is the expression
in~\eqref{eq:ehz-main} evaluated at $(\sigma, \beta)$.
Dividing by~$4T^2$:
\begin{equation}\label{eq:action-sum}
  \sum_{1 \leq j < i \leq |S|}
    \beta_{\sigma(i)}\, \beta_{\sigma(j)}\,
    \omega_0(n_{\sigma(j)},\, n_{\sigma(i)})
  = \frac{1}{2T}.
\end{equation}

\emph{Step~6: Conclude.}
Steps~4--5 establish a bijection between simple Reeb orbits
$(S, \sigma, \tau)$ with $\tau_k > 0$ and pairs
$(\sigma, \beta)$ with $\beta \in M(K)$,
$\beta_{\sigma(k)} > 0$ for all $k = 1, \ldots, |S|$.
Under this bijection, the action is
$A(\gamma) = T = \frac{1}{2} \bigl[
  \text{double sum in~\eqref{eq:action-sum}}
\bigr]^{-1}$.
Since $T$ decreases as the double sum increases:
\[
  c_{\mathrm{EHZ}}(K)
  = \min_{\sigma, \beta} T
  = \frac{1}{2} \left[
    \max_{\substack{\sigma \in S_F \\ \beta \in M(K)}}
    \sum_{1 \leq j < i \leq |S|}
      \beta_{\sigma(i)}\, \beta_{\sigma(j)}\,
      \omega_0(n_{\sigma(j)},\, n_{\sigma(i)})
  \right]^{-1}.
  \qedhere
\]
\end{proof}

% Jörn: text approved (692c271) — from "We now verify" to \end{algorithm*}
We now verify that the algorithm in
Section~\ref{subsec:general-case-algorithm} correctly implements this
optimization. We restate the algorithm for convenience.

\begin{algorithm*}[Algorithm~\ref{alg:ehz}, restated]
Computes~\eqref{eq:ehz-main}.

\smallskip\noindent
\emph{Input:} $(n_1, h_1), \ldots, (n_F, h_F)$.
\quad
\emph{Output:} $c_{\mathrm{EHZ}}(K)$.

\smallskip\noindent
For each nonempty subset $S \subseteq \{1, \ldots, F\}$
and each permutation $\sigma$ of $S$
(identified up to cyclic shifts):

\begin{enumerate}
\item \textbf{Define}
  the normal matrix $N \in \mathbb{R}^{|S| \times 4}$,
  action matrix $H \in \mathbb{R}^{|S| \times |S|}$,
  and height vector $\eta \in \mathbb{R}^{|S|}$:
  \begin{align*}
    N_{id} &= n_{\sigma(i),d},
      & i &= 1, \ldots, |S|, \quad d = 1, \ldots, 4, \\
    H_{ij} &= \begin{cases}
      \omega_0(n_{\sigma(i)}, n_{\sigma(j)}) & \text{if } i < j, \\
      0 & \text{if } i = j, \\
      \omega_0(n_{\sigma(j)}, n_{\sigma(i)}) & \text{if } i > j,
    \end{cases} \\
    \eta_i &= h_{\sigma(i)},
      & i &= 1, \ldots, |S|.
  \end{align*}

\item \textbf{Solve.}
  \begin{equation}\tag{\ref{eq:linear-system}}
    \begin{pmatrix} N^\top & 0 & 0 \\ H & -N & -\eta \\ \eta^\top & 0 & 0 \end{pmatrix}
    \begin{pmatrix} \beta \\ \lambda \\ \nu \end{pmatrix}
    = \begin{pmatrix} 0 \\ 0 \\ 1 \end{pmatrix},
    \qquad
    \beta \in \mathbb{R}^{|S|}, \quad \lambda \in \mathbb{R}^4, \quad \nu \in \mathbb{R}.
  \end{equation}

\item \textbf{Filter.}
  Discard $(S, \sigma)$ if the system has no solution, or
  if $\beta_i \leq 0$ for any~$i$.
  If the system has multiple solutions, choose any solution
  with $\beta_i > 0$ for all~$i$, if one exists.

\item \textbf{Evaluate.}
  Record $\beta^\top H \beta$
  (well-defined by Lemma~\ref{lem:well-defined}).
\end{enumerate}

\noindent
Return
$c_{\mathrm{EHZ}}(K) = \bigl(\max_{S,\sigma} \beta^\top H \beta\bigr)^{-1}$.
\end{algorithm*}

\medskip

\begin{lemma}[Double sum as quadratic form]\label{lem:H-quadratic}
% Jörn: text approved (b7a8bc0) — from \begin{lemma} to \end{proof}
% QC: for i > j, H_{ij} = omega_0(n_{sigma(i)}, n_{sigma(j)}),
%   which matches the summand in the double sum. Expanding
%   beta^T H beta using symmetry + zero diagonal gives the sum.
% Downstream: connects the double sum in the shoelace derivation
%   to the matrix form beta^T H beta. Gradient is nabla(beta^T H beta) = 2H beta.
\[
  \sum_{1 \leq j < i \leq |S|}
    \beta_{\sigma(i)}\, \beta_{\sigma(j)}\,
    \omega_0(n_{\sigma(j)},\, n_{\sigma(i)})
  = \tfrac{1}{2}\, \beta^\top H\, \beta.
\]
In particular,
$\nabla_{\beta}
  (\beta^\top H\, \beta) = 2H \beta$.
\end{lemma}

\begin{proof}
For $i > j$,
$H_{ij} = \omega_0(n_{\sigma(j)}, n_{\sigma(i)})$
by definition of~$H$.
Hence the double sum equals
$\sum_{j < i} \beta_{\sigma(i)}\, \beta_{\sigma(j)}\, H_{ij}$.
Since $H$ is symmetric with zero diagonal,
$\sum_{j < i} \beta_{\sigma(i)} \beta_{\sigma(j)} H_{ij}
= \tfrac{1}{2}\, \beta^\top H\, \beta$.
\end{proof}

\begin{lemma}[KKT conditions and the algorithm's linear system]%
\label{lem:kkt}
% Jörn: text approved (b7a8bc0) — from \begin{lemma} to \end{proof}
% QC: connects the abstract optimization (max beta^T H beta on M(K))
%   to the algorithm's concrete linear system (eq:linear-system).
%   With the new H sign convention, no sign flip is needed:
%   lambda = mu, nu = nu' directly.
For fixed $(S, \sigma)$, the linear
system~\eqref{eq:linear-system} is the Lagrange system for
critical points of
$\beta^\top H\, \beta$ on the constraint
manifold
$\{N^\top \beta = 0,\; \eta^\top \beta = 1\}$.
\end{lemma}

\begin{proof}
The gradient is
$\nabla_{\beta}
  (\beta^\top H \beta) = 2H \beta$
(Lemma~\ref{lem:H-quadratic}).
The Lagrange conditions
$\nabla(\beta^\top H \beta)
= N\lambda + \eta\nu$
give
$2H\beta = N\lambda + \eta\nu$.
Absorbing the factor~$2$ into the multipliers
($\lambda' = \lambda/2$, $\nu' = \nu/2$):
\[
  H \beta - N \lambda' - \eta \nu' = 0,
\]
which, together with $N^\top \beta = 0$ and
$\eta^\top \beta = 1$, is exactly
system~\eqref{eq:linear-system}.
\end{proof}

\medskip

\begin{lemma*}[Lemma~\ref{lem:well-defined}, restated]
If system~\eqref{eq:linear-system} has multiple solutions
$(\beta, \lambda, \nu)$, they all yield the same value
$\beta^\top H \beta$.
\end{lemma*}

\begin{proof}[Proof of Lemma~\ref{lem:well-defined}]
% Jörn: text approved (0bbc647) — from \begin{lemma*} to \end{proof}
Write the three block rows of system~\eqref{eq:linear-system} as
\[
  \text{(I)}\ N^\top \beta = 0, \qquad
  \text{(II)}\ H\beta = N\lambda + \eta\nu, \qquad
  \text{(III)}\ \eta^\top \beta = 1.
\]
For any solution, left-multiply~(II) by~$\beta^\top$:
\[
  \beta^\top H\beta
  = \beta^\top(N\lambda + \eta\nu)
  = \underbrace{(N^\top\beta)^\top}_{=\,0\text{ by (I)}}\lambda
    + \underbrace{(\eta^\top\beta)}_{=\,1\text{ by (III)}}\nu
  = \nu.
\]
Now let $(\delta\beta, \delta\lambda, \delta\nu)$ be any element
of the null space, so that
\[
  \text{(I$'$)}\ N^\top\delta\beta = 0, \qquad
  \text{(II$'$)}\ H\,\delta\beta = N\,\delta\lambda + \eta\,\delta\nu, \qquad
  \text{(III$'$)}\ \eta^\top\delta\beta = 0.
\]
Left-multiply~(II) by~$\delta\beta^\top$:
\[
  \delta\beta^\top H\beta
  = \underbrace{(N^\top\delta\beta)^\top}_{=\,0\text{ by (I$'$)}}\lambda
    + \underbrace{(\eta^\top\delta\beta)}_{=\,0\text{ by (III$'$)}}\nu
  = 0.
\]
Left-multiply~(II$'$) by~$\delta\beta^\top$:
\[
  \delta\beta^\top H\,\delta\beta
  = \underbrace{(N^\top\delta\beta)^\top}_{0}\delta\lambda
    + \underbrace{(\eta^\top\delta\beta)}_{0}\delta\nu
  = 0.
\]
Therefore
\[
  (\beta + \delta\beta)^\top H(\beta + \delta\beta)
  = \beta^\top H\beta
    + 2\,\delta\beta^\top H\beta
    + \delta\beta^\top H\,\delta\beta
  = \nu. \qedhere
\]
\end{proof}

% Jörn: math approved (8818fb8) — from \begin{lemma} to end of \begin{remark}

\begin{lemma}[Dismissal of redundant systems]%
\label{lem:rank-deficiency-dismissal}
Let $(S, \sigma)$ be a pair for which the linear
system~\eqref{eq:linear-system} has a solution
$(\beta_0, \lambda_0, \nu_0)$ with $\beta_0 \geq 0$,
and suppose the null space of the system matrix contains an element
$(\delta\beta, \delta\lambda, \delta\nu)$ with
$\delta\beta \neq 0$.

Then there exist $k \in S$ and a solution
$(\beta^*, \lambda^*, \nu^*)$ with
\begin{enumerate}
\item $\beta^*_k = 0$,
\item $\beta^*_j \geq 0$ for all $j \in S$,
\item $\beta^{*\top} H\, \beta^* = \beta_0^\top H\, \beta_0$.
\end{enumerate}
The restriction $\beta^*|_{S'}$ with $S' = S \setminus \{k\}$
lies in the constraint set $M(K)$ restricted to~$S'$,
so the algorithm for~$(S', \sigma|_{S'})$ maximises
$\beta^\top H'\, \beta$ over at least the value~$\nu_0$.
Consequently, the algorithm may discard any pair $(S, \sigma)$
whose system has such a null-space direction
without missing the capacity-achieving pair.
\end{lemma}

\begin{proof}
If $\beta_0$ already has $\beta_{0,k} = 0$ for some $k \in S$,
take $\beta^* = \beta_0$: conclusions (1)--(3) hold immediately
(with~(3) trivial), and the argument of Step~5 below gives the
restriction to~$S'$. So assume $\beta_{0,i} > 0$ for all~$i$.

Since $(\delta\beta, \delta\lambda, \delta\nu)$ is in the null space:
\[
  \text{(I$'$)}\ N^\top \delta\beta = 0, \qquad
  \text{(II$'$)}\ H\,\delta\beta
    = N\,\delta\lambda + \eta\,\delta\nu, \qquad
  \text{(III$'$)}\ \eta^\top \delta\beta = 0.
\]

\emph{Step~1: $\delta\beta$ has mixed signs.}
Expanding~(III$'$):
$\sum_{i \in S} \eta_i\, \delta\beta_i = 0$
where $\eta_i = h_{\sigma(i)} > 0$ are the heights.
Since $\delta\beta \neq 0$, at least one component is
nonzero, and the vanishing weighted sum with positive weights
forces $\delta\beta$ to have at least one strictly positive
and at least one strictly negative component.

\emph{Step~2: Walk along the null direction.}
Define $\beta(\alpha) = \beta_0 + \alpha\,\delta\beta$.
For every~$\alpha$, the triple
$(\beta(\alpha),\, \lambda_0 + \alpha\,\delta\lambda,\,
  \nu_0 + \alpha\,\delta\nu)$
satisfies~\eqref{eq:linear-system} (linearity of the system).
By Lemma~\ref{lem:well-defined}
(all solutions yield the same value of
$\beta^\top H\, \beta$),
$\beta(\alpha)^\top H\, \beta(\alpha) = \nu_0 + \alpha\,\delta\nu$.
On the other hand, expanding the quadratic form and using
$\delta\beta^\top H\, \beta_0 = 0$ and
$\delta\beta^\top H\, \delta\beta = 0$
(both from the proof of Lemma~\ref{lem:well-defined}):
\[
  \beta(\alpha)^\top H\, \beta(\alpha)
  = \beta_0^\top H\, \beta_0
    + 2\alpha\,\delta\beta^\top H\,\beta_0
    + \alpha^2\,\delta\beta^\top H\,\delta\beta
  = \nu_0.
\]
Comparing the two expressions: $\delta\nu = 0$, so the
multiplier~$\nu$ is the same for every solution.

\emph{Step~3: Identify the boundary point.}
Since $\beta_{0,i} > 0$ for all~$i$,
$\beta(\alpha)$ remains strictly positive for all
$|\alpha|$ sufficiently small.
For each $i$ with $\delta\beta_i < 0$, the component
$\beta_i(\alpha)$ reaches zero at
$\alpha_i = \beta_{0,i} / |\delta\beta_i|$.
Define
\[
  \alpha^*
  = \min\bigl\{\,
      \alpha_i : i \in S,\; \delta\beta_i < 0
    \bigr\},
\]
and let $k$ be an index achieving this minimum.
Such an index exists because $\delta\beta$ has at least one
negative component (Step~1).
Note $\alpha^* > 0$ since $\beta_{0,i} > 0$ and
$|\delta\beta_i| > 0$ for each candidate~$i$.

\emph{Step~4: Verify non-negativity.}
Set $\beta^* = \beta(\alpha^*)$.
Then $\beta^*_k = 0$ by construction.
For $j \neq k$:
\begin{itemize}
\item If $\delta\beta_j \geq 0$:
  $\beta^*_j = \beta_{0,j} + \alpha^*\,\delta\beta_j
    \geq \beta_{0,j} > 0$.
\item If $\delta\beta_j < 0$:
  $\alpha^* \leq \alpha_j = \beta_{0,j}/|\delta\beta_j|$
  (since $\alpha^*$ is the minimum), so
  $\beta^*_j
    = \beta_{0,j} - \alpha^*\,|\delta\beta_j|
    \geq \beta_{0,j} - \alpha_j\,|\delta\beta_j| = 0$.
\end{itemize}
Hence $\beta^*_j \geq 0$ for all $j \in S$.

\emph{Step~5: The smaller pair achieves the same value.}
Since $\beta^*_k = 0$, the quadratic form
$\beta^{*\top} H\, \beta^* = \nu_0$ (Step~2)
receives no contribution from index~$k$.
Let $H'$ denote the action matrix for
$(S', \sigma|_{S'})$, where $S' = S \setminus \{k\}$.
The sub-matrix of~$H$ at rows and columns~$S'$
equals~$H'$ (both are defined by the same normals
$n_{\sigma(j)}$, $j \in S'$), so
$(\beta^*|_{S'})^\top H'\, (\beta^*|_{S'}) = \nu_0$.

The restriction $\beta^*|_{S'}$ satisfies:
$\sum_{j \in S'} \beta^*_j\, n_{\sigma(j)} = 0$
(from $N^\top \beta^* = 0$ and $\beta^*_k = 0$),
$\sum_{j \in S'} \beta^*_j\, h_{\sigma(j)} = 1$
(from $\eta^\top \beta^* = 1$ and $\beta^*_k = 0$),
and $\beta^*_j \geq 0$ for $j \in S'$ (Step~4).
Hence $\beta^*|_{S'} \in M(K)$ restricted to~$S'$.
The algorithm evaluates
$\max\, \beta^\top H'\, \beta$ over this constraint set
and therefore finds value $\geq \nu_0$.
\end{proof}

\begin{remark}[Geometric characterisation of the null space]%
\label{rem:nullspace-geometry}
The null space of the system matrix always contains the subspace
$\{(0, \delta\lambda, 0) : N\,\delta\lambda = 0\}$, which has
dimension $4 - \operatorname{rank}(N)$.
In particular, a null-space element with $\delta\beta = 0$
exists if and only if the visited normals
$\{n_{\sigma(i)} : i \in S\}$ do not span~$\mathbb{R}^4$.
(Proof: if $\delta\beta = 0$, the null-space equations reduce to
$N\,\delta\lambda + \eta\,\delta\nu = 0$; multiplying by any
$\beta_0 \geq 0$ with $\eta^\top\!\beta_0 = 1$ gives $\delta\nu = 0$,
leaving $N\,\delta\lambda = 0$.)
The algorithm does not need this characterisation: it suffices
to inspect the null-space vectors from the SVD and check whether
$\delta\beta \neq 0$.
\end{remark}

For the numerical treatment of near-singular systems
(where the null space is only approximate),
see Appendix~\ref{app:near-singular}.

\begin{theorem}[Algorithm correctness]\label{thm:algorithm-correct}
% Jörn: text approved (0bbc647) — from \begin{theorem} to \end{proof}
% QC: the proof only uses completeness (the maximizer is found),
%   not soundness (every candidate is valid). This is intentional:
%   non-maximizing critical points need not correspond to
%   Reeb orbits (see Remark below).
% Downstream: the algorithm is proven correct. Any faithful
%   implementation computes c_EHZ.
The algorithm in Section~\ref{subsec:general-case-algorithm} correctly computes
$c_{\mathrm{EHZ}}(K)$.
\end{theorem}

\begin{proof}
By Theorem~\ref{thm:main},
$c_{\mathrm{EHZ}}(K)
= [\max \beta^\top H\, \beta]^{-1}$,
where the maximum is over all $\sigma \in S_F$ and
$\beta \in M(K)$.

\emph{Step~1: Existence.}
$M(K)$ is compact (closed and bounded in~$\mathbb{R}^F$),
$\beta^\top H\, \beta$ is continuous,
and the set of cyclic orderings~$\sigma$
is finite. Hence the maximum is attained at
some~$(\sigma^*, \beta^*)$ with value
$V^* = \beta^{*\top} H\, \beta^* > 0$
(since $c_{\mathrm{EHZ}}(K) > 0$).

\emph{Step~2: The maximizer is a candidate.}
Define $S = \{i : \beta^*_i > 0\}$.
Since $\beta^*$ maximizes
$\beta^\top H\, \beta$
over all of~$M(K)$, it also
maximizes it over the face
$\{\beta \in M(K) : \beta_i = 0 \;\forall\, i \notin S\}$.
As $\beta^*_i > 0$ for all $i \in S$, this is an interior
maximum on the face, so the Lagrange conditions hold:
$\beta^*$ satisfies the linear
system~\eqref{eq:linear-system} for~$(S, \sigma^*)$
(Lemma~\ref{lem:kkt}).
The filter passes: $\beta^*_i > 0$ for all $i \in S$ by
definition, so the algorithm chooses some solution with
$\beta_i > 0$.
By Lemma~\ref{lem:well-defined}, the recorded value
$\beta^\top H\, \beta = V^*$ regardless of which
solution is chosen.

% FIXME(hbox): 11.92pt — "$\max_{S,\sigma} \beta^\top H\, \beta = V^*$" and "$(V^*)^{-1} = c_{\mathrm{EHZ}}(K)$" on one line; try \linebreak before "and the algorithm" or split into two sentences
\emph{Step~3: The algorithm returns it.}
For any other candidate~$(S', \sigma')$,
$\beta'^\top H\, \beta' \leq V^*$
by maximality. Hence
$\max_{S,\sigma} \beta^\top H\, \beta = V^*$,
and the algorithm returns
$(V^*)^{-1} = c_{\mathrm{EHZ}}(K)$.
\end{proof}

\begin{remark}\label{rem:non-maximizer-warning}
% Jörn: text approved (b7a8bc0)
% QC: important for agents/readers who might assume all critical
%   points correspond to Reeb orbits.
The proof uses only that the global maximizer
of~$\beta^\top H\, \beta$
corresponds to a minimum-action Reeb orbit
(Theorem~\ref{thm:main}, Step~6).
Non-maximizing critical points with $\beta > 0$
need not correspond to Reeb orbits:
the correspondence between $(\sigma, \beta)$ and simple orbits
is established only at the global maximum.
The algorithm may record $\beta^\top H\, \beta$ at such
spurious critical points, but they cannot affect the
returned value since they are not the maximum
(Step~3 above).
\end{remark}

